\documentclass[../cours_1re.tex]{subfiles}
\begin{document}
\setcounter{chapter}{17}
\chapter{Droites, projections et cercles}
\label{ch:18}
\begin{objectif}
Passer d'une figure à des équations ; construire un projeté orthogonal et reconnaître un cercle.
Prérequis : équations de droites de seconde et produit scalaire.
\end{objectif}
\section{Équation cartésienne et vecteur normal}
\begin{propriete}
Si $(a,b)\ne(0,0)$, l'ensemble $d$ d'équation $ax+by+c=0$
est une droite de vecteur normal $\vec n=(a,b)$ et de vecteur directeur
$(-b,a)$.
\end{propriete}
\begin{demonstration}
Il existe un point $A=(x_A,y_A)$ satisfaisant l'équation :
si $a\ne0$, prendre $(-c/a,0)$ ; sinon prendre $(0,-c/b)$.
Un point $M$ vérifie l'équation si et seulement si
$a(x-x_A)+b(y-y_A)=0$, c'est-à-dire $\vec n\cdot\overrightarrow{AM}=0$.
Les vecteurs perpendiculaires à $\vec n$ sont précisément les multiples de
$(-b,a)$, ce qui se vérifie en résolvant cette équation linéaire.
\end{demonstration}
Une droite passant par $A$ et normale à $(a,b)$ a donc pour équation
$a(x-x_A)+b(y-y_A)=0$. Multiplier tous les coefficients par un même
nombre non nul ne change pas la droite.
Deux droites sont perpendiculaires si et seulement si leurs vecteurs
normaux sont orthogonaux.

\section{Projection orthogonale et distance à une droite}
\begin{theoreme}
Pour $M=(x_0,y_0)$, le projeté orthogonal $H$ sur $d:ax+by+c=0$ est
\[
H=M-\frac{ax_0+by_0+c}{a^2+b^2}(a,b).
\]
La distance de $M$ à $d$ vaut
$\displaystyle \frac{|ax_0+by_0+c|}{\sqrt{a^2+b^2}}$.
\end{theoreme}
\begin{demonstration}
La droite passant par $M$ et normale à $d$ est paramétrée par $M-t(a,b)$.
Substituer ses coordonnées dans l'équation de $d$ donne l'unique valeur
$t=(ax_0+by_0+c)/(a^2+b^2)$. On a donc $H\in d$ et $MH\perp d$.
La norme de $\overrightarrow{MH}=-t(a,b)$ donne la distance.
Pour tout $P\in d$, Pythagore donne $MP^2=MH^2+HP^2\ge MH^2$ :
le projeté est le point de $d$ le plus proche de $M$, et il est unique.
\end{demonstration}
\begin{exemple}
Pour $d:3x+4y-5=0$ et $M=(3,4)$, le numérateur vaut $20$.
$H=(3,4)-\frac45(3,4)=(3/5,4/5)$ ; on vérifie l'équation de $d$.
La distance vaut $20/5=4$.
\end{exemple}

\section{Reconnaître un cercle}
Le cercle de centre $\Omega=(a,b)$ et de rayon $r>0$ a pour équation
\[
(x-a)^2+(y-b)^2=r^2.
\]
C'est exactement la condition $\Omega M=r$, par la formule de la distance.
Réciproquement, compléter les carrés dans
$x^2+y^2+ux+vy+w=0$ donne
\[
(x+u/2)^2+(y+v/2)^2=(u^2+v^2)/4-w.
\]
Le membre de droite positif définit un cercle ; nul, il définit un point ;
négatif, il n'y a aucun point réel.
\begin{exemple}
$x^2+y^2-4x+6y+9=0$ devient $(x-2)^2+(y+3)^2=4$ :
centre $(2,-3)$, rayon $2$.
\end{exemple}
La tangente en $T$ au cercle est perpendiculaire à $\Omega T$.
Elle a pour équation
$(x_T-a)(x-x_T)+(y_T-b)(y-y_T)=0$.
Pourquoi une seule intersection ? Pour un autre point $P$ de cette droite,
$\Omega P^2=\Omega T^2+TP^2>r^2$. Elle touche donc seulement le point $T$.

\section{Intersections et choix du calcul}
Pour intersecter un cercle avec une droite non verticale, remplacer $y$
par $mx+p$ dans son équation : on obtient un second degré en $x$.
Pour une droite verticale $x=k$, remplacer directement $x$ par $k$.
Le discriminant correspond à deux points, une tangence ou aucun point.
On peut aussi comparer la distance du centre à la droite avec le rayon.
\begin{exemple}
$x^2+y^2=25$ et $y=3$ donnent $x^2=16$, donc $(-4,3)$ et $(4,3)$.
Avec $y=5$, il reste seulement $(0,5)$ ; avec $y=6$, aucun point.
\end{exemple}

\section{Une parabole comme lieu géométrique}
\approfondissement
Fixons $p>0$, un foyer $F=(0,p)$ et une directrice $d:y=-p$.
Les points équidistants de $F$ et de $d$ vérifient
\[
\sqrt{x^2+(y-p)^2}=|y+p|
\iff x^2+(y-p)^2=(y+p)^2
\iff y=\frac{x^2}{4p}.
\]
L'élévation au carré est réversible parce que les deux membres sont positifs
ou nuls. On retrouve une parabole de sommet l'origine et d'axe vertical.
La droite horizontale $y=k$ coupe cette parabole en deux points si $k>0$,
un si $k=0$, aucun si $k<0$ ; une verticale la coupe toujours une fois.

\begin{finchapitre}
\section{Utilité et prolongements}
Les équations permettent de vérifier une construction et de calculer
des distances minimales. La normalité intervient aussi dans les tangentes
à des courbes et dans les contraintes d'optimisation.

En licence, le gradient généralisera le vecteur normal à des courbes et surfaces.
Les projections sur des sous-espaces, étudiées au chapitre 27, reprendront
exactement la preuve par Pythagore de la meilleure approximation.
\begin{retenir}
Une normale décrit une droite. Le projeté vérifie deux conditions :
appartenir à la droite et donner une direction perpendiculaire. Compléter les carrés classe les lieux.
\end{retenir}
\end{finchapitre}
\end{document}
